\documentclass[11pt]{article}
\usepackage[T1]{fontenc}
\usepackage[a4paper,margin=1in]{geometry}
\usepackage{amsmath,amssymb,amsthm}
\usepackage[hidelinks]{hyperref}
\usepackage{microtype}

\newtheorem*{questions}{Source Questions 3.2}
\newtheorem*{answer}{Later answer (Corollary 2.3)}

\setlength{\parindent}{0pt}
\setlength{\parskip}{0.6em}

\title{Literature Status: Isomorphic Invariance and\\
Reflexivity for the Weak Maximizing Property}
\author{Run \texttt{fa\_banach\_001} --- lane 10}
\date{11 August 2026}

\begin{document}
\maketitle

\begin{center}
\fbox{\parbox{0.9\textwidth}{
\textbf{Status: literature already answered.}
Questions~3.2 of arXiv:1906.02162 ask whether the weak maximizing
property (WMP) is invariant under isomorphism and whether WMP is equivalent
to reflexivity.  The separate later paper arXiv:2103.10288 explicitly answers
both questions negatively in Corollary~2.3.
}}
\end{center}

\section{Original questions}

For Banach spaces $E,F$, the pair $(E,F)$ has the WMP if every bounded
operator $T:E\to F$ possessing a non-weakly-null maximizing sequence attains
its norm.  Aron, Garc\'ia, Pellegrino, and Teixeira ask on arXiv PDF
page~9~\cite{source}:

\begin{questions}
Are there infinite-dimensional spaces $E,F_1,F_2$ such that $(E,F_1)$ has
the WMP but $(E,F_2)$ does not?  Is WMP preserved when the domain or range is
replaced by an isomorphic copy?  Is a Banach space reflexive exactly when it
has the WMP?
\end{questions}

The source had already proved that WMP implies reflexivity, so the final
question asks whether the converse holds.

\section{Separate later answer}

Dantas, Jung, and Mart\'inez-Cervantes state in the introduction of
arXiv:2103.10288 that they answer Questions~3.2.  Their Corollary~2.3, on
arXiv PDF pages~4--5, says:

\begin{answer}
There are reflexive Banach spaces $E_1,E_2,E_3$, all isomorphic to $\ell_2$,
such that $(\ell_2,E_1)$, $(E_2,\ell_2)$, and $E_3$ fail the WMP.
\end{answer}

The proof applies their Theorem~2.1 to a non-norm-attaining operator on
$\ell_2$.  Concretely, it uses the renormed finite sums
$\ell_2\oplus_q\mathbb R$ for $q<2$ and
$\ell_2\oplus_p\mathbb R$ for $p>2$, including $p=\infty$.

Since $(\ell_2,\ell_2)$ has the WMP, while $E_1$ is isomorphic to $\ell_2$
and $(\ell_2,E_1)$ fails it, WMP is not preserved by an isomorphic change of
the range.  Taking $E=\ell_2$, $F_1=\ell_2$, and $F_2=E_1$ answers the first
existence question.  The space $E_3$ is reflexive but fails WMP, so
reflexivity does not imply WMP.  The examples involving $E_2$ also show
failure under an isomorphic change of the domain.

Thus every component of Questions~3.2 receives a negative answer.  The later
paper explicitly identifies these consequences; this packet records a
literature match and claims no new theorem.

\section{Scope}

The separate question in the source about WMP for $(L_p,L_q)$ is not answered
in full by this identification.  The later paper proves failure when $p>2$
or $q<2$ and leaves the standard-norm range $1<p\leq2\leq q<\infty$ open.

\begin{thebibliography}{9}
\bibitem{source}
R. M. Aron, D. Garc\'ia, D. Pellegrino, and E. V. Teixeira,
\emph{Reflexivity and non-weakly null maximizing sequences},
Proc. Amer. Math. Soc. \textbf{148} (2020), 741--750;
arXiv:1906.02162.

\bibitem{answerpaper}
S. Dantas, M. Jung, and G. Mart\'inez-Cervantes,
\emph{Some remarks on the weak maximizing property},
arXiv:2103.10288.
\end{thebibliography}

\end{document}
